\chapter{Stationary Stochastic Processes}

\section{Strict Stationary Processes}

\begin{definition}[Strict-Sense Stationary Process]\index{strict-sense stationary process}
  A \emph{strict-sense stationary process} (SSS) is a stochastic process \(\{X_t\}_{t \in T}\) such that for any constant \(h\), for any positive integer \(n\), for any \(t_1, t_2, \dots, t_n \in T\), and for any \(t_1 + h, t_2 + h, \dots, t_n + h \in T\),
  \[
    (X_{t_1}, X_{t_2}, \dots, X_{t_n})
  \]
  and
  \[
    (X_{t_1 + h}, X_{t_2 + h}, \dots, X_{t_n + h})
  \]
  have the same \(n\)th-order probability distribution function, i.e.,
  \[
    F(x_1, x_2, \dots, x_n; t_1, t_2, \dots, t_n) = F(x_1, x_2, \dots, x_n; t_1 + h, t_2 + h, \dots, t_n + h),
  \]
  for \(x_1, x_2, \dots, x_n \in \mathbb{R}\).
\end{definition}

\begin{theorem}
  Suppose that the process \(\{X_t\}_{t \in T}\) is strict-sense stationary and of second moments. Then,
  \begin{enumerate}
    \item for all \(t \in T\), \(\mu_X(t) = \mathbb{E}[X_t]\) is a constant.
    \item \(R_X(s, t) = \mathbb{E}[X_s X_t]\) depends only on the time difference \(t - s\) for all \(s, t \in T\).
  \end{enumerate}
\end{theorem}

Furthermore,
\begin{align*}
  C_X(t, t + \tau) &= \mathbb{E}[(X_t - \mu_X(t))(X_{t + \tau} - \mu_X(t + \tau))] \\
  &= \mathbb{E}[X_t X_{t + \tau}] - \mathbb{E}[X_t \mu_X] - \mathbb{E}[\mu_X X_{t + \tau}] + \mathbb{E}[\mu_X^2] \\
  &= R_X(\tau) - \mu_X^2 = C_X(\tau)
\end{align*}
depends only on \(\tau\) and
\[
  \sigma_X^2(t) = C_X(0) = R_X(0) - \mu_X^2
\]
is a constant for all \(t \in T\).

\begin{example}
  Supppose that \(\{X_n\}_{n \geq 0}\) are IID random sequences, each of which has a uniform distribution in the interval \((0, 1)\). Is \(\{X_n\}_{n \geq 0}\) a strict-sense stationary process? Determine the value \(\mathbb{E}[X_n]\) and \(\mathbb{E}[X_n X_m], n, m = 0, 1, 2, \dots\).
\end{example}
\begin{solution}
  Suppose that the distribution function of \(X_n\) is \(F(x)\). Then for all \(k, h \in \mathbb{N}\) and \(0 < n_1 < n_2 < \dots < n_k\), the joint distribution function of \((X_{n_1}, X_{n_2}, \dots, X_{n_k})\) and \((X_{n_1 + h}, X_{n_2 + h}, \dots, X_{n_k + h})\) are both
  \begin{align*}
    F(x_1, x_2, \dots, x_k) &= P(X_{n_1} \leq x_1, X_{n_2} \leq x_2, \dots, X_{n_k} \leq x_k) \\
    &= P(X_{n_1} \leq x_1) P(X_{n_2} \leq x_2) \cdots P(X_{n_k} \leq x_k) \\
    &= F(x_1) F(x_2) \cdots F(x_k).
  \end{align*}
  Thus, \(\{X_n\}_{n \geq 0}\) is a strict-sense stationary process.

  The expected value of \(X_n\) is
  \[
    \mathbb{E}[X_n] = \frac{1}{2},
  \]
  since \(X_n\) is uniformly distributed on the interval \((0, 1)\).

  For \(n = m\),
  \[
    \mathbb{E}[X_n X_m] = \mathbb{E}[X_n^2] = \int_0^1 x^2 \, \mathrm{d}x = \frac{1}{3}.
  \]
  
  For \(n \neq m\),
  \[
    \mathbb{E}[X_n X_m] = \mathbb{E}[X_n] \mathbb{E}[X_m] = \left(\frac{1}{2}\right)^2 = \frac{1}{4},
  \]
  due to the independence of \(X_n\) and \(X_m\).
\end{solution}

\section{Wide Stationary Processes}

\begin{definition}
  A process \(\{X_t\}_{t \in T}\) is called a \emph{wide-sense stationary process} (WSS) if it is of second order and satisfies the following conditions:
  \begin{enumerate}
    \item for all \(t \in T\), \(\mu_X(t) = \mathbb{E}[X_t]\) is a constant.
    \item for all \(t, t + \tau \in T\), \(R_X(t, t + \tau) = R_X(\tau)\) depends only on the time difference \(\tau\).
  \end{enumerate}
  In particular, a random sequence \(\{X_n\}_{n \geq 0}\) is called a stationary random sequence (or time sequence) in the wide sense if it is of second order and satisfies the following conditions:
  \begin{enumerate}
    \item for all \(n \geq 0\), \(\mathbb{E}[X_n]\) is a constant.
    \item for all \(m, n > \geq 0\), \(R_X(n, n + m) = R_X(m)\) depends only on the time difference \(m\).
  \end{enumerate}
\end{definition}